\documentclass[aspectratio=169]{beamer}

\usepackage[utf8]{inputenc}
\usepackage[T1]{fontenc}
\usepackage{amsmath,amssymb}
\usepackage{listings}
\usepackage{xcolor}
\usepackage{array}
\usepackage{booktabs}
\usepackage{hyperref}

% ── Colour palette ─────────────────────────────────────────────────────────────
\definecolor{orangetheme}{RGB}{220,95,0}
\definecolor{pyblue}    {RGB}{42,115,205}
\definecolor{pygreen}   {RGB}{0,128,0}
\definecolor{pygray}    {RGB}{110,110,110}
\definecolor{pyred}     {RGB}{180,0,0}
\definecolor{pybg}      {RGB}{248,248,252}
\definecolor{linkblue}  {RGB}{0,100,200}

% ── Beamer theme ───────────────────────────────────────────────────────────────
\usetheme{default}
\setbeamercolor{title}         {fg=orangetheme}
\setbeamercolor{frametitle}    {fg=orangetheme}
\setbeamercolor{structure}     {fg=orangetheme}
\setbeamercolor{itemize item}  {fg=orangetheme}
\setbeamercolor{itemize subitem}{fg=orangetheme}
\setbeamerfont{title}    {size=\Large,series=\bfseries}
\setbeamerfont{frametitle}{size=\large,series=\bfseries}
\setbeamertemplate{navigation symbols}{}
\setbeamertemplate{footline}[frame number]

% ── Python listings style ──────────────────────────────────────────────────────
\lstdefinestyle{python}{
  language        = Python,
  basicstyle      = \footnotesize\ttfamily,
  keywordstyle    = \color{pyblue}\bfseries,
  stringstyle     = \color{pyred},
  commentstyle    = \color{pygreen}\itshape,
  backgroundcolor = \color{pybg},
  frame           = single,
  framerule       = 0.4pt,
  rulecolor       = \color{gray!40},
  showstringspaces= false,
  breaklines      = true,
  xleftmargin     = 4pt,
  xrightmargin    = 4pt,
  aboveskip       = 4pt,
  belowskip       = 4pt,
  morekeywords    = {True,False,None},
}

% ── Meta ───────────────────────────────────────────────────────────────────────
\title{Predicates and Quantifiers}
\subtitle{Discrete Mathematics --- Lesson 4}
\author{Notes By Pr. El Hadiq Zouhair}
\date{}

% ══════════════════════════════════════════════════════════════════════════════
\begin{document}

% ── Title slide ────────────────────────────────────────────────────────────────
\begin{frame}
  \titlepage
\end{frame}

% ── 1. Overview ────────────────────────────────────────────────────────────────
\begin{frame}{Overview}
  \begin{itemize}
    \item Predicate logic expands the range of propositions that are possible.
    \item A particular form of programming is possible using predicate logic:
          \textcolor{linkblue}{logic programming}.
    \item Universal truths can be differentiated from special cases.
    \item You will learn about predicates and quantifiers.
  \end{itemize}
\end{frame}

% ── 2. Predicates ──────────────────────────────────────────────────────────────
\begin{frame}[fragile]{Predicates}
  Many propositions have variables that determine their value.
  \begin{itemize}
    \item The melting point of $x$ element is $y$ Kelvin: $M(x,\,y)$
  \end{itemize}
  Popular predicates include equality predicates:
  $<,\;>,\;\leq\;(\texttt{<=}),\;\geq\;(\texttt{>=}),\;=\;(\texttt{==}),\;\neq\;(\texttt{!=})$

  \begin{lstlisting}[style=python]
from sympy import symbols, Implies

x = symbols('x')
# Predicate: (x + 2 > 4) => (x + 3 < 4)
pred = Implies(x + 2 > 4, x + 3 < 4)

# Evaluate for x = 2 and x = 3
results = [pred.subs(x, v) for v in [2, 3]]
print(results)  # [True, False]
  \end{lstlisting}
\end{frame}

% ── 3. Example: Pythagorean triple ─────────────────────────────────────────────
\begin{frame}[fragile]{Example: Pythagorean Triple}
  Consider the proposition $p\colon\; x^{2} + y^{2} = z^{2}$

  \begin{lstlisting}[style=python]
# Define the predicate as a Python lambda
p = lambda x, y, z: x**2 + y**2 == z**2

# What is the truth value for x=5, y=12, z=13?
print(p(5, 12, 13))     # True

# What is the truth value for x=160, y=241, z=290?
print(p(160, 241, 290)) # False
  \end{lstlisting}

  \begin{itemize}
    \item To learn more about this equation, read about
          \textcolor{linkblue}{Pythagorean triples}.
  \end{itemize}
\end{frame}

% ── 4. Universal Quantifier ────────────────────────────────────────────────────
\begin{frame}[fragile]{Universal Quantifier $\boldsymbol{\forall}$}
  Universal quantification: ``$P(x)$ is true for \emph{all} values of~$x$.''

  \[
    \forall x\;(x^{2} \geq 0)
  \]

  \begin{lstlisting}[style=python]
from z3 import Real, ForAll, Implies, prove

x = Real('x')
# For all real x: x^2 >= 0
prove(ForAll([x], x**2 >= 0))                  # proved
  \end{lstlisting}

  \[
    \forall x \geq 2\;(x \geq 0)
  \]

  \begin{lstlisting}[style=python]
# For all x >= 2: x >= 0  (encode domain as implication)
prove(ForAll([x], Implies(x >= 2, x >= 0)))    # proved
  \end{lstlisting}
\end{frame}

% ── 5. Existential Quantifier ──────────────────────────────────────────────────
\begin{frame}[fragile]{Existential Quantifier $\boldsymbol{\exists}$}
  Existential quantification: ``$P(x)$ is true for \emph{at least one} value of~$x$.''

  \[
    \exists x\;(x = -x)
  \]

  \begin{lstlisting}[style=python]
from z3 import Real, Reals, And, Exists, Solver, Goal, Tactic, sat

x = Real('x')
s = Solver()
s.add(x == -x)
print(s.check())   # sat
print(s.model())   # [x = 0]
  \end{lstlisting}

  This is unique: $\exists!\,x$. \quad
  $\exists\,x,y\;\bigl(x^{2}+ay^{2}\leq 1 \;\wedge\; x-y \geq 2\bigr)$

  \begin{lstlisting}[style=python]
x, y, a = Reals('x y a')
# Quantifier elimination: remove x, y to find condition on a
g = Goal()
g.add(Exists([x, y], And(x**2 + a*y**2 <= 1, x - y >= 2)))
print(Tactic('qe')(g))   # [[a <= 1/3]]
  \end{lstlisting}
\end{frame}

% ── 6. Example: Region inclusion ───────────────────────────────────────────────
\begin{frame}[fragile]{Example: Region Inclusion}
  Test whether region $R_{1}$ is contained in region $R_{2}$:

  \begin{lstlisting}[style=python]
from z3 import Reals, And, Not, Solver

x, y = Reals('x y')
R1_z3 = (2*x)**2 + y**2 + 2*x*y <= 1
R2_z3 = x**2 + y**2 <= 2

# Exists (x,y): R1 and not R2?
s = Solver()
s.add(And(R1_z3, Not(R2_z3)))
print(s.check())   # unsat  =>  R1 subset R2
  \end{lstlisting}

  \begin{lstlisting}[style=python]
import numpy as np, matplotlib.pyplot as plt
X, Y = np.meshgrid(np.linspace(-2,2,500), np.linspace(-2,2,500))
R1 = (2*X)**2 + Y**2 + 2*X*Y <= 1
R2 = X**2 + Y**2 <= 2
plt.contourf(X, Y, R2, levels=[.5,1.5], colors=['lightblue'], alpha=.5)
plt.contourf(X, Y, R1, levels=[.5,1.5], colors=['orange'],    alpha=.7)
plt.axis('equal'); plt.show()
  \end{lstlisting}
\end{frame}

% ── 7. Equivalences ────────────────────────────────────────────────────────────
\begin{frame}{Equivalences}
  Equivalences are also used for quantifiers
  (verified with \texttt{sympy.simplify} or \texttt{reduce\_inequalities}):

  \medskip
  \begin{center}
    \renewcommand{\arraystretch}{1.45}
    \begin{tabular}{|l|c|}
      \hline
      \textbf{Law} & \textbf{Equivalence} \\
      \hline
      Conjunction &
        $\forall x\bigl(P(x)\wedge Q(x)\bigr)
         \;\Leftrightarrow\;
         \forall x\,P(x)\wedge\forall x\,Q(x)$ \\
      Disjunction &
        $\exists x\bigl(P(x)\vee Q(x)\bigr)
         \;\Leftrightarrow\;
         \exists x\,P(x)\vee\exists x\,Q(x)$ \\
      De Morgan's $(\forall)$ &
        $\neg\forall x\,P(x)
         \;\Leftrightarrow\;
         \exists x\,\neg P(x)$ \\
      De Morgan's $(\exists)$ &
        $\neg\exists x\,P(x)
         \;\Leftrightarrow\;
         \forall x\,\neg P(x)$ \\
      \hline
    \end{tabular}
  \end{center}

  \smallskip
  The negation can be pushed all the way through an expression:
  \begin{align*}
    &\neg\exists v\;\forall w\;\exists x\;\exists y\;\forall z\;P(v,w,x,y,z)\\
    \Leftrightarrow\;&\forall v\;\neg\forall w\;\exists x\;\exists y\;\forall z\;P(v,w,x,y,z)\\
    \Leftrightarrow\;&\forall v\;\exists w\;\neg\exists x\;\exists y\;\forall z\;P(v,w,x,y,z)\\
    \Leftrightarrow\;&\forall v\;\exists w\;\forall x\;\forall y\;\neg\forall z\;P(v,w,x,y,z)\\
    \Leftrightarrow\;&\forall v\;\exists w\;\forall x\;\forall y\;\exists z\;\neg P(v,w,x,y,z)
  \end{align*}
\end{frame}

% ── 8. Nested Quantifiers ──────────────────────────────────────────────────────
\begin{frame}{Nested Quantifiers}
  How to interpret nested quantifiers:

  \medskip
  \begin{center}
    \renewcommand{\arraystretch}{1.45}
    \begin{tabular}{|c|l|}
      \hline
      $\exists x\;\exists y\;P(x,y)$ & At least one pair exists. \\
      $\forall x\;\exists y\;P(x,y)$ & A pair exists for every $x$. \\
      $\exists x\;\forall y\;P(x,y)$ & A specific $x$ works for any $y$. \\
      $\forall x\;\forall y\;P(x,y)$ & All pairs work! \\
      \hline
    \end{tabular}
  \end{center}

  \medskip
  Everyone in this course likes a branch of mathematics:
  \[
    \forall\,s \in \mathrm{students}\;\;
    \exists\,b \in \mathrm{mathematics}\;\;
    \bigl(\mathrm{Likes}(s,\,b)\bigr)
  \]
\end{frame}

% ── 9. Examples: Translating statements ────────────────────────────────────────
\begin{frame}{Examples: Translating Statements to Logic}
  Translate the following statements into logic:

  \medskip
  \begin{itemize}
    \item \emph{No one is perfect.}
    \[
      \neg\,\exists x\;\mathrm{Perfect}(x)
    \]

    \item \emph{Not all propositions are contradictions.}
    \[
      \neg\,\forall x\;\mathrm{Contradiction}(x)
    \]

    \item \emph{In my second-year class, only one person knew Java, SQL and Python.}
    \[
      \exists!\,x\;\bigl(\mathrm{Java}(x)\;\wedge\;\mathrm{SQL}(x)\;\wedge\;\mathrm{Python}(x)\bigr)
    \]
  \end{itemize}
\end{frame}

% ── 10. Summary ────────────────────────────────────────────────────────────────
\begin{frame}{Summary}
  \begin{itemize}
    \item Predicates are propositional functions with parameters.
    \item There are three quantifiers:
    \begin{itemize}
      \item \textbf{Universal} \quad(\texttt{z3.prove(ForAll([x],\;...))},\; $\forall$)
      \item \textbf{Existence} \quad(\texttt{z3.Solver().check()},\; $\exists$)
      \item \textbf{Uniqueness} \quad($\exists!$)
    \end{itemize}
    \item Quantifier elimination via \texttt{z3.Tactic('qe')} mirrors Wolfram's \texttt{Reduce}.
    \item So far, propositions were taken as is.
    \item The next lesson allows for analysis of propositions through
          \textbf{rules of inference}.
  \end{itemize}
\end{frame}

\end{document}
